% !TEX program = xelatex
% =============================================================================
%  “华为杯”中国研究生数学建模竞赛 —— 官方 LaTeX 模板（GMCMthesis）
%
%  来源：第二十三届华为杯latex模板（数模加油站版）
%       相对官方原版的唯一修改见 gmcmthesis.cls 顶部的「Linux/macOS 适配」注释。
%
%  编译链（照抄官方 makefiles.sh，四趟，缺一不可）：
%      xelatex main.tex
%      bibtex  main          ← 漏掉这趟就没有 main.bbl，\cite 全变 [?]
%      xelatex main.tex
%      xelatex main.tex
%    或直接执行 ./compile.sh
%
%  字体：fontset=windows 需要 SimSun / SimHei / KaiTi / FangSong 四种真品字体。
%        若本机没有，改成 fontset=mac（用 Songti SC / Heiti SC / Kaiti SC）
%        或直接删掉 fontset= 选项走 ctex 默认（fandol，字形与官方有出入）。
% =============================================================================
\documentclass[bwprint,fontset=windows]{gmcmthesis}

\hypersetup{hidelinks}
\usepackage[framemethod=TikZ]{mdframed}
\usepackage{subfig}
\usepackage{siunitx}
\usepackage{colortbl}
\usepackage[noend]{algpseudocode}
\usepackage{algorithmicx,algorithm}
\floatname{algorithm}{算法}
\renewcommand{\algorithmicrequire}{\textbf{输入:}}
\renewcommand{\algorithmicensure}{\textbf{输出:}}

\newcommand{\red}[1]{\textcolor{red}{#1}}
\newcommand{\blue}[1]{\textcolor{blue}{#1}}

% 图片搜索路径。\includegraphics 的路径基准是【主文件 main.tex 所在的目录】，
% 不是写出该行的 section 文件——所以在 sections/*.tex 里也只写裸文件名即可，
% 由下面这几条前缀兜底（paper/figures、项目根/figures、再上一层/figures）。
\graphicspath{{figures/}{../figures/}{../../figures/}}

% -----------------------------------------------------------------------------
% 代码清单样式：逐项照抄 cls 里 \lstnewenvironment{Python} 的外观（彩色关键字 +
% 浅灰底 + 细框 + 文件名标题条），另补四项附录用 \lstinputlisting 时才需要的设置。
% 详见本模板说明，不要随手改成「纯黑无高亮」——那与模板外观不一致。
% -----------------------------------------------------------------------------
\lstset{
  frame=single, framesep=4pt, framerule=0.5pt,
  rulecolor=\color{gray!50},
  backgroundcolor=\color{gray!5},
  xleftmargin=10pt, xrightmargin=10pt,
  keywordstyle=\color{pKeyword},
  stringstyle=\color{pString},
  commentstyle=\color{pComment},
  tabsize=4,
  captionpos=t,
  basicstyle=\linespread{1}\small\ttfamily,   % \linespread 必须在字号切换之前
  upquote=true,
  columns=flexible,
  extendedchars=true,
  breaklines=true,
  breakatwhitespace=false,
  showstringspaces=false,
  postbreak=\mbox{\textcolor{gray}{$\hookrightarrow$}\space},
}

% -----------------------------------------------------------------------------
% 封面信息（今年封面为西安交通大学 logo + 官方 title 图，不要替换 figures/ 下的两份 PDF）
% -----------------------------------------------------------------------------
\title{【论文标题】}
\baominghao{\hspace{6em} 【参赛队号】}
\schoolname{\hspace{5.8em} 【学校名称】}
\membera{\hspace{6em} 【队员A】}
\memberb{\hspace{6em} 【队员B】}
\memberc{\hspace{6em} 【队员C】}

\begin{document}

% 封面
\maketitle

% 摘要页：正文写在 abstract.tex，关键词用 \keywords（会自动进摘要页底部）
\begin{abstract}
% 摘要：由 main.tex 的 \begin{abstract} 环境 \input 进来。
% 中文摘要必须写全：问题概述 + 每个子问题用了什么方法 + 每个子问题的精确数值结果 + 结论。
% 不要写「本文建立了模型，效果良好」这类没有数值的空话。
% 关键词由 main.tex 里的 \keywords{...} 给出（不要写在本文件里）。

本文针对【问题背景一句话】问题，围绕【子问题一】【子问题二】【子问题三】展开研究。

针对问题一，【数据与口径】。建立【模型名称】，【求解方法】。结果表明，【关键数值结果】。

针对问题二，【数据与口径】。建立【模型名称】，【求解方法】。结果表明，【关键数值结果】。

针对问题三，【数据与口径】。建立【模型名称】，【求解方法】。结果表明，【关键数值结果】。

最后，对【关键参数】进行了灵敏度分析，【结论】。

\keywords{关键词一 \quad 关键词二 \quad 关键词三}
\end{abstract}

\pagestyle{plain}

% 目录：必须用 \maketoc（模板的壳，含 \phantomsection + PDF 书签「目录」）
\maketoc
\newpage

\section{问题重述}

\subsection{问题背景}

【背景一段。只做客观交代，不复述题目原文。】

\subsection{问题的提出}

围绕【核心对象】，本文依次解决以下问题：

\textbf{问题一：}【一句话，含给定条件】

\textbf{问题二：}【一句话，含给定条件】

\textbf{问题三：}【一句话，含给定条件】
    % 问题重述
\section{问题分析}

\subsection{总体分析}

【整体思路：数据 -> 各子问题模型 -> 结果检验。说明子问题之间如何递进、
前一问的哪个中间量是后一问的输入。数据流程图（如有）也放这一节。】

% 技术路线图：把 fig_roadmap.pdf 放进 figures/ 后取消注释。
% 路径只写裸文件名——\includegraphics 的基准是 main.tex 所在目录，不是本文件。
% \begin{figure}[H]
%   \centering
%   \includegraphics[width=0.92\textwidth]{fig_roadmap.pdf}
%   \caption{本文技术路线}
%   \label{fig:roadmap}
% \end{figure}

\subsection{问题一分析}

【问题一的目标、可用数据、难点、拟采用的模型类型。】

\subsection{问题二分析}

\subsection{问题三分析}
       % 问题分析
\section{模型假设}

\begin{enumerate}
  \item 【假设一，并说明为什么合理】
  \item 【假设二】
  \item 【假设三】
\end{enumerate}
    % 模型假设
\section{符号说明}

\begin{table}[H]
  \centering
  \caption{符号说明}
  \label{tab:symbols}
  \begin{tabular}{cll}
    \toprule
    符号 & 含义 & 单位 \\
    \midrule
    $x_i$      & 【含义】 & 【单位】 \\
    $\beta_j$  & 【含义】 & --- \\
    \bottomrule
  \end{tabular}
\end{table}

% 表/式/图 的编号都是「节号.序号」（cls 里 \thetable/\theequation/\thefigure 均如此），
% 例如第 4 节的第一个表格就是「表 4.1」。写 ref 时用 \ref{tab:...} 让 LaTeX 自己算。
        % 符号说明
\section{问题一：【短标题】}

\subsection{模型建立}

【模型假设、变量定义、目标函数与约束。公式编号由 cls 自动生成，
形如 (5.1)——cls 里 \theequation 就是「节号.序号」。】

\begin{equation}
  \label{eq:p1_obj}
  \min \sum_{i=1}^{n} c_i x_i
\end{equation}

\subsection{求解方法}

【算法步骤、停止条件。算法的「输入/输出」标签已由模板设成中文，直接用即可。】

\begin{algorithm}[H]
  \caption{【算法名】}
  \begin{algorithmic}[1]
    \Require 【输入】
    \Ensure  【输出】
    \State 【步骤一】
    \State 【步骤二】
  \end{algorithmic}
\end{algorithm}

\subsection{求解结果}

【关键数值必须与本次求解的实际结果一致，不得编造、估算或另行取整。】

\begin{table}[H]
  \centering
  \caption{【结果表题】}
  \label{tab:p1_result}
  \begin{tabular}{cccc}
    \toprule
    【列一】 & 【列二】 & 【列三】 & 【列四】 \\
    \midrule
    【数据】 & 【数据】 & 【数据】 & 【数据】 \\
    \bottomrule
  \end{tabular}
\end{table}

% 数据图同样只写裸文件名：
% \begin{figure}[H]
%   \centering
%   \includegraphics[width=0.8\textwidth]{fig_q1_error.pdf}
%   \caption{【图题】}
%   \label{fig:p1_error}
% \end{figure}
       % 问题一
\section{问题二：【短标题】}

\subsection{模型建立}

【可引用问题一的结论，例如「问题一给出的 $x_i$ 在本节作为约束右端项」。】

\subsection{求解方法}

\subsection{求解结果}

【关键数值 + 结果图表。】
       % 问题二
\section{问题三：【短标题】}

\subsection{模型建立}

\subsection{求解方法}

\subsection{求解结果}

\subsection{结果讨论}

【与问题一/二的对比，说明新增条件的边际效果。】
       % 问题三
\section{灵敏度分析}

\subsection{【关键参数一】的影响}

\subsection{【关键参数二】的影响}

% 灵敏度曲线用数据图（3coding-visual 生成），不要用 drawio 画。
    % 灵敏度分析
\section{模型评价与推广}

\subsection{模型优点}

\begin{enumerate}
  \item 【优点一，要具体到方法，不要写「模型具有科学性」这类空话】
  \item 【优点二】
\end{enumerate}

\subsection{模型不足}

\begin{enumerate}
  \item 【不足一，并说明它对结论的影响有多大】
\end{enumerate}

\subsection{模型推广}

【这套建模框架还能用到哪些场景。】
     % 模型评价与推广

% 顺序：正文 → 参考文献 → 附录（参考文献必须排在附录之前）
% 参考文献：走官方模板的 BibTeX 机制（gmcm.bst = gbt7714-unsrt.bst）。
%
% 文献条目写在 reference.bib，正文用 \cite{键} 引用。
% 编译链必须是 xelatex → bibtex → xelatex → xelatex（见 ../compile.sh）。
%
% 两个静默失败点：
%   ① bibtex 只收录「被 \cite 过」的条目——没引用过的条目不会出现在 PDF 里；
%   ② 漏跑 bibtex 时 \cite 渲染成 [?]，而 xelatex 不报任何错。
%
% 位置：main.tex 里 % 参考文献：走官方模板的 BibTeX 机制（gmcm.bst = gbt7714-unsrt.bst）。
%
% 文献条目写在 reference.bib，正文用 \cite{键} 引用。
% 编译链必须是 xelatex → bibtex → xelatex → xelatex（见 ../compile.sh）。
%
% 两个静默失败点：
%   ① bibtex 只收录「被 \cite 过」的条目——没引用过的条目不会出现在 PDF 里；
%   ② 漏跑 bibtex 时 \cite 渲染成 [?]，而 xelatex 不报任何错。
%
% 位置：main.tex 里 % 参考文献：走官方模板的 BibTeX 机制（gmcm.bst = gbt7714-unsrt.bst）。
%
% 文献条目写在 reference.bib，正文用 \cite{键} 引用。
% 编译链必须是 xelatex → bibtex → xelatex → xelatex（见 ../compile.sh）。
%
% 两个静默失败点：
%   ① bibtex 只收录「被 \cite 过」的条目——没引用过的条目不会出现在 PDF 里；
%   ② 漏跑 bibtex 时 \cite 渲染成 [?]，而 xelatex 不报任何错。
%
% 位置：main.tex 里 \input{references} 排在附录之前，顺序不能反。
\bibliographystyle{gmcm}
\bibliography{reference}
 排在附录之前，顺序不能反。
\bibliographystyle{gmcm}
\bibliography{reference}
 排在附录之前，顺序不能反。
\bibliographystyle{gmcm}
\bibliography{reference}


\begin{appendices}
% 本文件被 main.tex 的 \begin{appendices} ... \end{appendices} 包住，
% 所以这里只写 \section，不要自己写 appendices 环境。
% \section 会被渲染成「附录 A」，子小节是 A.1、A.2……

\section{源程序}

以下为本文全部计算程序的清单。所有代码在 Python 3.12 环境下可直接运行，
按问题一至问题三的顺序执行即可完整复现正文中的全部数值结果。

\renewcommand{\thesubsection}{A\thinskip.\thinskip\arabic{subsection}}

% ── 放什么，不放什么 ──────────────────────────────────────────────
% 只放【参与算出论文里数值】的计算程序：utils + 各子问题求解程序。
% 【不放】生成流程图 / 示意图 / drawio XML 的脚本（如 make_diagrams.py）——
% 那类代码正文是 drawio 样式串，单行可达 160+ 列（版心只容 68 列），
% 在附录里每行折成两三行，对评委是噪音又不产生任何数值结论。

% 文件名写在代码框的标题条上（等价于模板的 \begin{Python}{文件名}），
% 子节标题只写说明，不要再重复文件名。
% 路径基准：\lstinputlisting 相对【主文件 main.tex 所在的目录】解析，不是相对本文件。
\subsection{数据加载与共享工具}
\lstinputlisting[language=Python, title=utils.py]{../code/utils.py}

\subsection{问题一：【短标题】}
\lstinputlisting[language=Python, title=problem1.py]{../code/problem1.py}

\subsection{问题二：【短标题】}
\lstinputlisting[language=Python, title=problem2.py]{../code/problem2.py}

\subsection{问题三：【短标题】}
\lstinputlisting[language=Python, title=problem3.py]{../code/problem3.py}

\end{appendices}

\end{document}
